\documentclass[../../main.tex]{subfiles}
\begin{document}
\section{Frobenius 定理}\label{section:ch19-frobenius}

在例\cref{example:ch19-分布与积分流形}中，除最后一个之外我们定义的分布都具有这样的性质：过每一点都有一条积分流形。而且在局部，这些子流形都像 $\mathbb{R}^n$ 中平行的仿射子空间那样优美地"拼合"在一起。给定一个秩-$k$ 分布 $D\subseteq TM$，若 $M$ 上的一个光滑坐标图 $(U,\varphi)$ 满足 $\varphi(U)$ 是 $\mathbb{R}^n$ 中的一个立方体，且在 $U$ 的点处 $D$ 由前 $k$ 个坐标向量场 $\partial/\partial x^1,\dots,\partial/\partial x^k$ 张成（见图\cref{figure:ch19-flat-chart}），则称 $(U,\varphi)$ 是 $D$ 的一个\emph{平坦坐标图}。在任意这样的坐标图中，形如 $x^{k+1}=c^{k+1},\dots,x^n=c^n$（$c^{k+1},\dots,c^n$ 为常数）的每个切片都是 $D$ 的积分流形。这是积分流形最理想的局部情形。

\begin{figure}[ht]
\centering
\includegraphics[width=0.8\textwidth]{fig-19-2.png}
\caption{分布的平坦图}\label{figure:ch19-flat-chart}
\end{figure}

若在 $M$ 的每一点的一个邻域内都存在 $D$ 的一个平坦坐标图，则称分布 $D\subseteq TM$ 是\emph{完全可积的}。显然，每个完全可积分布都是可积的，因而是对合的。概括地说：
\begin{align*}
\text{完全可积}\Longrightarrow\text{可积}\Longrightarrow\text{对合}.
\end{align*}

下一个定理是本章的主要结果，确实也是光滑流形理论中最核心的定理之一。它说上述蕴含关系实际上都是等价关系：
\begin{align*}
\text{完全可积}\Longleftrightarrow\text{可积}\Longleftrightarrow\text{对合}.
\end{align*}

\begin{theorem}[Frobenius 定理]\label{theorem:ch19-frobenius}
每个对合分布都是完全可积的。
\end{theorem}

\begin{proof}
交换向量场的标准型定理（定理\cref{theorem:ch09-46}）表明，局部地由一组独立的光滑交换向量场张成的任意分布都是完全可积的，因为该定理所保证存在的坐标图是平坦的（在必要时缩小定义域使像为一个立方体）。因此，只需证明每个对合分布都局部地由一组独立的光滑交换向量场张成。

设 $D$ 是 $n$ 维流形 $M$ 上一个秩为 $k$ 的对合分布，$p\in M$。由于完全可积性是一个局部问题，通过过渡到 $p$ 的一个光滑坐标邻域，可把 $M$ 替换为 $\mathbb{R}^n$ 的一个开子集 $U$，并为 $D$ 选取一组光滑局部标架 $X_1,\dots,X_k$。必要时重新排列坐标，可设 $D_p$ 与 $T_p\mathbb{R}^n$ 中由 $\partial/\partial x^{k+1}|_p,\dots,\partial/\partial x^n|_p$ 张成的子空间互补。

令 $\pi\colon\mathbb{R}^n\to\mathbb{R}^k$ 为到前 $k$ 个坐标的投影，$\pi(x^1,\dots,x^n)=(x^1,\dots,x^k)$（见图\cref{figure:ch19-frobenius-proof}）。这诱导出一个光滑丛同态 $\mathrm{d}\pi\colon T\mathbb{R}^n\to T\mathbb{R}^k$，它可写成
\begin{align*}
\mathrm{d}\pi\left(\sum_{i=1}^n v^i\frac{\partial}{\partial x^i}\bigg|_q\right)=\sum_{i=1}^k v^i\frac{\partial}{\partial x^i}\bigg|_{\pi(q)}.
\end{align*}
（注意右端的求和只对 $i=1,\dots,k$ 进行。）由于 $\mathrm{d}\pi|_D$ 是包含 $D\hookrightarrow TU$ 与 $\mathrm{d}\pi$ 的复合，故它是一个光滑丛同态。于是，关于标架 $\{X_i|_q\}$ 与 $\{\partial/\partial x^i|_{\pi(q)}\}$，$\mathrm{d}\pi|_{D_q}$ 的矩阵元在 $D$ 的每一根纤维上都是非退化的。因此它有一个光滑逆，从而定义了 $D$ 的光滑局部截面 $V_1,\dots,V_k$，使得 $V_i$ 与 $\partial/\partial x^i$ 是 $\pi$-相关的。由于 $\mathrm{d}\pi[V_i,V_j]=[\mathrm{d}\pi V_i,\mathrm{d}\pi V_j]=[\partial/\partial x^i,\partial/\partial x^j]=0$，而 $D$ 的对合性蕴含 $[V_i,V_j]$ 取值于 $D$ 中，且 $\mathrm{d}\pi$ 在 $D$ 的每一根纤维上都是单射，故这蕴含对每个 $q$ 都有 $[V_i,V_j]_q=0$，从而完成证明。
\end{proof}

\begin{figure}[ht]
\centering
\includegraphics[width=0.8\textwidth]{fig-19-3.png}
\caption{Frobenius定理的证明}\label{figure:ch19-frobenius-proof}
\end{figure}

为了后面处理超定偏微分方程时使用，我们注意到由该证明得到的一个简单推论。

\begin{corollary}\label{corollary:ch19-平坦切片}
设 $M$ 是一个光滑流形，$D$ 是 $M$ 上一个对合的秩-$k$ 分布，$S\subseteq M$ 是一个余维 $k$ 的嵌入子流形。若 $p\in S$ 满足 $T_pS$ 与 $D_p$ 互补，则存在一个以 $p$ 为中心的 $D$ 的平坦坐标图 $(U,\theta^i)$，使得 $S\cap U$ 是切片 $\theta^1=\cdots=\theta^k=0$。
\end{corollary}

\begin{proof}
该定理的证明表明，局部地 $D$ 由 $k$ 个交换向量场张成，于是本推论由定理\cref{theorem:ch09-46} 得到。
\end{proof}

通常，在 Frobenius 定理的证明中嵌入了一种寻找积分流形的技巧。其想法是利用坐标投影找到张成同一分布的交换向量场，再用例\cref{example:ch09-47} 的技巧找到平坦坐标图。下面是一个例子。

\begin{example}\label{example:ch19-frobenius-example}
设 $D\subseteq T\mathbb{R}^3$ 是由下列向量场张成的分布：
\begin{align*}
X=x\frac{\partial}{\partial x}+\frac{\partial}{\partial y}+x(y+1)\frac{\partial}{\partial z},\qquad
Y=\frac{\partial}{\partial x}+y\frac{\partial}{\partial z}.
\end{align*}
例\cref{example:ch08-27} 的计算表明
\begin{align*}
[X,Y]=-\frac{\partial}{\partial x}-y\frac{\partial}{\partial z}=-Y,
\end{align*}
故 $D$ 对合。我们来求原点邻域内的一个平坦坐标图。由于 $D$ 与 $\partial/\partial z$ 的张成空间互补，坐标投影 $\pi\colon\mathbb{R}^3\to\mathbb{R}^2$，$\pi(x,y,z)=(x,y)$，在每一点 $(x,y,z)\in\mathbb{R}^3$ 诱导出同构 $\mathrm{d}\pi|_{D_{(x,y,z)}}\colon D_{(x,y,z)}\to T_{(x,y)}\mathbb{R}^2$。Frobenius 定理的证明表明，若能找到 $D$ 的光滑局部截面 $V,W$，它们分别与 $\partial/\partial x$、$\partial/\partial y$ 是 $\pi$-相关的，则它们就是张成 $D$ 的交换向量场。容易验证，$V,W$ 具有此性质当且仅当它们取值于 $D$ 中，且形如
\begin{align*}
V=\frac{\partial}{\partial x}+u(x,y,z)\frac{\partial}{\partial z},\qquad
W=\frac{\partial}{\partial y}+v(x,y,z)\frac{\partial}{\partial z},
\end{align*}
其中 $u,v$ 是某些光滑实值函数。一点线性代数表明，向量场
\begin{align*}
V=Y=\frac{\partial}{\partial x}+y\frac{\partial}{\partial z},\qquad
W=X-xY=\frac{\partial}{\partial y}+x\frac{\partial}{\partial z}
\end{align*}
满足要求。这两个向量场的流可通过求解两组常微分方程方便地求得。省去细节，我们求得 $V$ 的流为
\begin{align}
\alpha_t(x,y,z)=(x+t,y,z+ty),
\end{align}
$W$ 的流为
\begin{align}
\beta_t(x,y,z)=(x,y+t,z+tx).
\end{align}
于是按例\cref{example:ch09-47} 的步骤，可从 $z$ 轴出发、依次沿着这两个流向外流动，来定义坐标映射的逆 $\Phi$：
\begin{align*}
\Phi(u,v,w)=\alpha_u\circ\beta_v(0,0,w)=\alpha_u(0,v,w)=(u,v,w+uv).
\end{align*}
我们所求的平坦坐标由反演映射 $(x,y,z)=\Phi(u,v,w)=(u,v,w+uv)$ 得到：
\begin{align*}
(u,v,w)=\Phi^{-1}(x,y,z)=(x,y,z-xy).
\end{align*}
由此 $D$ 的积分流形是 $w(x,y,z)=z-xy$ 的水平集。（由于这里构造的平坦坐标图实际上是一个整体坐标图，这描述了所有的积分流形，而不仅是原点附近的那些。）
\end{example}

\begin{exercise}\label{exercise:ch19-15}
验证 $V$ 与 $W$ 的流分别由上面两式给出，并验证 $z-xy$ 的水平集是 $D$ 的积分流形。
\end{exercise}

下一个命题是 Frobenius 定理最重要的推论之一。

\begin{proposition}[积分流形的局部结构]\label{proposition:ch19-局部结构}
设 $D$ 是光滑流形 $M$ 上一个秩为 $k$ 的对合分布，$(U,\theta^i)$ 是 $D$ 的一个平坦坐标图。若 $H$ 是 $D$ 的任意积分流形，则 $H\cap U$ 是 $U$ 的平行 $k$ 维切片的可数多个互不相交开子集的并，其中每一个在 $H$ 中都是开的，且在 $M$ 中是嵌入的。
\end{proposition}

\begin{proof}
设 $H$ 是 $D$ 的积分流形。由于包含映射 $\iota\colon H\hookrightarrow M$ 连续，$H\cap U=\iota^{-1}(U)$ 在 $H$ 中是开的，因而由可数多个连通分支组成，其中每一个在 $H$ 中都是开的。

令 $V$ 为 $H\cap U$ 的任意一个连通分支（见图\cref{figure:ch19-local-structure}）。我们首先证明 $V$ 含于单个切片中。由于 $\mathrm{d}\theta^{k+1},\dots,\mathrm{d}\theta^n$ 是 $D$ 的局部定义形式，可知这些 $1$-形式拉回到 $V$ 上恒为零。因为 $V$ 连通，这蕴含 $\theta^{k+1},\dots,\theta^n$ 在 $V$ 上全为常数，故 $V$ 含于单个切片 $S$ 中。由于 $S$ 在 $M$ 中嵌入，由推论\cref{corollary:ch05-30}，包含映射 $V\hookrightarrow M$ 作为到 $S$ 的映射也是光滑的。于是包含映射 $V\hookrightarrow S$ 是相同维数流形之间的单射光滑浸入，因而是局部微分同胚、开映射，并且是到 $S$ 的一个开子集上的同胚。包含映射 $V\hookrightarrow M$ 是光滑嵌入 $V\hookrightarrow S\hookrightarrow M$ 的复合，故它是一个光滑嵌入。
\end{proof}

\begin{figure}[ht]
\centering
\includegraphics[width=0.8\textwidth]{fig-19-4.png}
\caption{积分流形的局部结构}\label{figure:ch19-local-structure}
\end{figure}

上述命题给出关于积分流形的一个重要结果，我们将在本章末尾研究 Lie 子群时使用。回忆，若每个像含于 $H$ 的光滑映射 $F\colon N\to M$ 作为从 $N$ 到 $H$ 的映射都是光滑的，则称光滑子流形 $H\subseteq M$ 在 $M$ 中是\emph{弱嵌入的}（见第 5 章）。

\begin{theorem}\label{theorem:ch19-弱嵌入}
对合分布的每个积分流形都是弱嵌入的。
\end{theorem}

\begin{proof}
设 $M$ 是光滑 $n$ 维流形，$H\subseteq M$ 是 $M$ 上一个对合的秩-$k$ 分布 $D$ 的积分流形，$F\colon N\to M$ 是满足 $F(N)\subseteq H$ 的光滑映射。任取 $p\in N$，令 $q=F(p)\in H$。设 $(y^1,\dots,y^n)$ 是 $D$ 在 $q$ 的邻域 $U$ 上的平坦坐标，$(x^i)$ 是 $N$ 在 $p$ 的一个连通邻域 $B$ 上的光滑坐标，使得 $F(B)\subseteq U$。把 $F$ 的坐标表示写成 $(y^1,\dots,y^n)=(F^1(x),\dots,F^n(x))$。$F(B)\subseteq H\cap U$ 这一事实意味着坐标函数 $F^{k+1},\dots,F^n$ 只取可数多个值。由于 $B$ 连通，介值定理蕴含这些坐标函数为常数，从而 $F(B)$ 含于单个切片 $S\subseteq U$ 中。由于 $S\cap H$ 是 $H$ 的一个在 $M$ 中嵌入的开子集，可知 $F|_B$ 作为从 $B$ 到 $S\cap H$ 的映射是光滑的，从而由复合，$F|_B\colon B\to(S\cap H)\hookrightarrow H$ 作为到 $H$ 的映射是光滑的。
\end{proof}
\end{document}
